% 20260522_Assingment_Probability_Statistics
\documentclass[dvipdfmx]{jsreport}
%preamble
\usepackage{soupreamble}

\title{確率及び統計学特論A 第2回課題}
\author{M1 6012 川村聡一郎}
\date{}


\begin{document}
\maketitle

\begin{tcolorbox}
	r.v. $X, Y$に対して
\begin{equation*}
	\mathcal{P}(X\le a)= \mathcal{P}(Y\le a),\ a\in \R
\end{equation*}
	ならば
\begin{equation*}
	\mathcal{P}(X\in A)= \mathcal{P}(Y\in A),\ A\in \mathcal{B}(\R)
\end{equation*}
	であることを示せ。
\end{tcolorbox}


\begin{souanswer} %答案はここに書く
	解答はここにかく
\end{souanswer}
\end{document}